% Источники: exam/materials/моделирование.pdf, разделы 47 и 18; exam/materials/преподаватель/2020/14-04-2020-Model-L9-1-14.04.20.pdf; exam/materials/прошлогодние ответы/Modelirovanie_1_-3.pdf, вопрос 18; GitHub HumsterProgrammer/iu9-notes/8sem/Моделирование/Моделирование_лекция-05_26-02-26.md; GitHub HumsterProgrammer/iu9-notes/8sem/Моделирование/Моделирование_лаб-02_16-03-26.md.
\section{Методика построения имитационной модели систем массового обслуживания.}

В имитационном моделировании участвуют несколько специалистов: специалисты предметной области, математики и IT-специалисты.

\subsection*{Методика построения}

\begin{enumerate}
  \item \textbf{Постановка задачи} в терминах предметной области: классификация сложной динамической системы (структурно-сложная, меняющая поведение во времени, гибридная и др.).

  \item \textbf{Построение структурной схемы} исследуемой системы с целью выделения компонентов и определения взаимодействия между ними.

  \item По схеме \textbf{уточняются известные входные данные}, параметры системы; определяются неизвестные промежуточные переменные, выходные данные, ограничения, обратная связь, управляющие воздействия, начальные и краевые условия. На этом этапе модель окончательно классифицируется как сложная, предварительно выбирается схема организации квазипараллелизма.

  \item Совместно со специалистом предметной области \textbf{строится концептуальная модель} исследуемой системы.

  \item \textbf{Преобразование концептуальной модели в имитационную}: выбор дискретно-событийной, процессной, транзактной или агрегатной схемы.

  \item \textbf{Уточнение значений} входных переменных и параметров системы или получение методики их оценивания.

  \item \textbf{Формализация модели}: описание потоков заявок, каналов обслуживания, накопителей, маршрутов, дисциплин очереди и правил отказа.

  \item \textbf{Отладка модели на ЦВМ.}

  \item \textbf{Счёт по программе.}

  \item \textbf{Проверка качества модели на практике}, оценка необходимости модификации.

  \item Совместно со специалистами предметной области и IT-специалистами \textbf{выявляются «узкие места»} модели. Для сокращения обсуждения уточнение модели начинается с этапа концептуального проектирования.

  \item \textbf{Проведение вычислительного эксперимента} с использованием построенной математической модели.
\end{enumerate}

\subsection*{Концептуальная модель СМО}

\textbf{СМО} — система, функционирование которой состоит в обработке заявок. Заявка — транзакт; очередь — канал ожидания заявок; обслуживающий аппарат — устройство обслуживания. Если источник входного потока расположен вне системы — СМО \textbf{открытая}; если заявка после обслуживания возвращается в поток — \textbf{закрытая}.

\textbf{Для СМО в концептуальной модели фиксируют:}
\begin{itemize}
  \item интенсивность и закон входного потока;
  \item число каналов, законы времени обслуживания и правила маршрутизации;
  \item ёмкость очередей, дисциплину обслуживания, приоритеты;
  \item показатели качества: время ожидания, длина очереди, загрузка каналов, вероятность отказа.
\end{itemize}

\subsection*{Построение модели СМО}

Для построения модели СМО:
\begin{enumerate}
  \item выполняется постановка задачи: что дано, что найти;
  \item строится схема: отмечаются заявки, обслуживающие устройства, каналы, распределения времени и условия окончания моделирования;
  \item по схеме записывается алгоритм; в рамках курса для этого используется GPSS.
\end{enumerate}

В лабораторных замечаниях отдельно фиксируется, что модель должна быть именно имитационной моделью СМО: нужна схема или рисунок, на котором видны очереди, обслуживающие устройства, входные и выходные потоки; далее по схеме раскрываются блоки, переходы, петли и соответствие GPSS-командам.

\textbf{Проверка качества} включает верификацию программы, валидацию модели по данным предметной области и статистическую обработку результатов серии прогонов.

\clearpage
